% Owner: cyj. Frozen model cyj.ndqp.scenario.v8; figures and values are
% traced to outputs/Q2/curated_manifest.json and the pinned upstream manifest.
\section{问题二：跨维度数据融合与广义标度律}

\subsection{数据角色与跨来源统一假设}

为同时刻画参数规模 $N$、训练 token 数 $D$、质量 $Q$ 和 17 维配比
$\boldsymbol p$，先固定证据角色。B1 的 1176 条 Pythia 轨迹拟合经典
$N$--$D$ 骨架；B7 的 450 条半合成记录只估计质量扩展。B2 为半合成
Cerebras 轨迹，B3 为 B1 风格插值轨迹，二者检查形状；B4/B5 只检查
同族规模变化的 Loss 方向。B6 包含在 B7 网格中，不能再作独立测试；
B8 与 B7 的质量方向有冲突，保持独立审计。B9 为大模型元数据，B10
的 Loss 为 estimated，只作外推压力参考。全文取 $N$ 为十亿参数、$D$
为十亿 token，Loss 为来源自身报告的交叉熵量。不同来源的 tokenizer、
验证语料和 Loss 聚合口径未证明一致，故没有统一绝对 Loss 的外部 RMSE。

问题一给出 A 侧质量评分 $Q_A$、17 域配比及 13 个目标域的相对 Loss
曲面；B7 的 \texttt{Q\_score} 是 B 侧原生半合成质量坐标，记作 $Q_B$。
现有附件既无同一次实验的 $(Q_A,Q_B)$ 成对观测，也无 A 侧逐目标 Loss
与 B1/B7 总体 \texttt{val\_loss} 的实测对齐。以下桥接均为明示的条件性
代理假设，不是经验标定。

\subsection{经典 $N$--$D$ 标度律建立}

以 B1 为唯一经典骨架来源，在 $N\in[0.070542,11.965825]$、
$D\in[0.134,299.893]$ 拟合
\begin{equation}\label{eq:cyj-b1}
L_0(N,D)=E+A N^{-\alpha}+B D^{-\beta}.
\end{equation}
估计对 8 个规模组等权，采用对数 Loss 残差的 Huber 目标、正参数约束
及固定种子的多起点搜索。以所述十亿单位为输入，$E,A,B$ 为 Loss
单位，$\alpha,\beta$ 无量纲。得到 $E=1.689838$、$A=0.353969$、
$B=1.240275$、$\alpha=0.339985$、$\beta=0.279892$。
全样本 RMSE 为 $1.4658\times10^{-4}$ Loss；留一规模组平均 RMSE 为
$1.4613\times10^{-4}$，各组最后 30\% token 检查点留出的 RMSE 为
$1.1600\times10^{-4}$。B1 的 $D$ 与步数乘固定每步 token 数的显示
舍入相容，\texttt{ppl} 与 $\exp(\texttt{val\_loss})$ 的显示舍入也全行
相容；附件缺少原始评估记录和生成脚本。这些极小残差只支持 B1 同源
条件重构，不代表真实跨来源预测有相同精度。

\subsection{数据质量扩展与模型选择}

冻结 B1 五参数，仅用 B7 估计三个质量系数：
\begin{align}\label{eq:cyj-b7}
L_B(N,D,Q_B)&=L_0(N,D)+(1-Q_B)G(N,D),\\
G(N,D)&=G_0+G_N\ln N+G_D\ln(D/100).
\end{align}
在 B7 固定 $(N,D)$ 的 45 个组中，质量斜率并不恒定；恒斜率近似
的组级 RMSE 为 $0.1078$。三系数在质量增益非负约束下最小二乘
估计为 $G_0=0.371703$、$G_N=-0.061115$、$G_D=-0.014961$。
四角检查得到 B7 矩形内 $G\geq0.193184$，故模型域内 $L_{Q_B}<0$。

\begin{table}[htbp]
\centering\small
\caption{最终模型的冻结参数；输入 $N,D$ 均以十亿单位计，$E,A,B,G_0,G_N,G_D$ 以 Loss 计，指数无量纲}
\label{tab:cyj-q2-parameters}
\begin{tabular}{p{0.15\textwidth}p{0.68\textwidth}}
\hline
拟合来源 & 估计参数\\ \hline
B1 经典骨架 & $E=1.689838,\ A=0.353969,\ B=1.240275,\ \alpha=0.339985,\ \beta=0.279892$\\
B7 质量扩展 & $G_0=0.371703,\ G_N=-0.061115,\ G_D=-0.014961$；B1 五参数保持固定\\ \hline
\end{tabular}
\end{table}

\begin{table}[htbp]
\centering\small
\caption{B7 半合成数据上的质量模型比较；RMSE 为 B7 原生 Loss，留出按 $N,D,Q_B$ 等级分折}
\label{tab:cyj-q2-quality}
\begin{tabular}{p{0.36\textwidth}rrrr}
\hline
模型 & 训练 & 留 $N$ & 留 $D$ & 留 $Q_B$\\ \hline
固定 B1 骨架＋3 质量参数 & 0.048540 & 0.048738 & 0.048751 & 0.048639\\
B7 联合 8 参数比较模型 & 0.048410 & 0.049764 & 0.049008 & 0.049352\\ \hline
\end{tabular}
\end{table}

三参数模型采用 24 折内层选择惩罚、外层留等级的嵌套验证；八参数
行是固定族留级比较，选择流程不同，不以微小差异宣称真实训练中的
优劣。选择固定 B1 骨架，是为了保持 B1 估计规模、B7 补充质量
的可追溯角色。B7 有 240 条记录的 $N,D$ 落在 B1 矩形外；它们参与
质量项估计，却不扩大最终预测范围。质量效应随规模变化的结果属于
半合成网格上的条件关系，不是质量干预的因果效应。

\subsection{质量与领域配比的跨来源条件融合}

设 $\mathcal M$ 为问题一中有 A1 质量映射的 6 个训练域（3 个 direct、
3 个 near-direct）。主方案只对已映射域计算均值：
\begin{equation}\label{eq:cyj-qa}
Q_A(\boldsymbol p)=\frac{\sum_{i\in\mathcal M}p_iq_{A,i}}
{\sum_{i\in\mathcal M}p_i},\qquad \sum_{i\in\mathcal M}p_i>0.
\end{equation}
无映射评分保留缺失；零覆盖的请求被拒绝。令 $q_{\min},q_{\max}$
为已映射域 A1 评分极值，主方案把该均值单调归一化为 B7 范围内
的代理坐标
\begin{equation}\label{eq:cyj-qproxy}
q=g(Q_A)=\operatorname{clip}_{[0.1,1]}
\left[0.1+0.9\frac{Q_A-q_{\min}}{q_{\max}-q_{\min}}\right].
\end{equation}
参考配比有 $Q_A^{\rm ref}=0.783104$、$q^{\rm ref}=0.435406$。
映射斜率的 0.5 倍和 1.5 倍另作敏感性情景；$q$ 始终是条件性
$Q_B$ 代理，不是观测到的 $Q_B$。

问题一冻结的 interaction surface 对 13 个目标域给出
$\widehat L_{A,k}(\boldsymbol p)$，令
\begin{align}\label{eq:cyj-rp}
R_k&=\widehat L_{A,k}(\boldsymbol p_{\rm ref})>0,\qquad
r_k(\boldsymbol p)=\frac{\widehat L_{A,k}(\boldsymbol p)-R_k}{R_k},\\
r_w(\boldsymbol p)&=\sum_{k=1}^{13}w_kr_k(\boldsymbol p),\qquad
w_k\geq0,\quad\sum_kw_k=1.
\end{align}
$w_k$ 是目标域分析权重，并非训练域配比。$r_w$ 在参考配比
为零；它是 A 侧 1M 条件下的相对 Loss 效应，向 B 侧总体 Loss
的迁移幅度不能由现有附件识别。

\subsection{最终广义标度律}

在上述明示的跨来源假设下，唯一条件主模型为
\begin{equation}\label{eq:cyj-v8}
\boxed{\widehat L(N,D,\boldsymbol p;\boldsymbol w)=
\left[L_0(N,D)+\{1-g(Q_A(\boldsymbol p))\}G(N,D)\right]
\exp\{r_w(\boldsymbol p)\}.}
\end{equation}
B1 拟合 $L_0$，B7 半合成数据拟合 $G$；$g$ 与指数乘子是
跨来源条件桥接。指数形式在参考配比乘子为 1、保持正 Loss，
但不证明实际 B 侧训练配比等于参考配比。$g$ 与 $r_w$ 均依赖
$\boldsymbol p$，两条路径的影响不能由附件完全分离。
式\eqref{eq:cyj-b7} 在 $Q_B=1$ 时退化为经典式\eqref{eq:cyj-b1}。
最终式在 $q=1$ 时仅质量项消失，配比修正仍可能存在；冻结参考
配比的代理值为 $0.435406$，不能再把 $q=1$ 与参考配比当作
同一个已观测情景。因此最终式不施加“任意配比下 $q=1$ 即退化”
的额外约束。

正式预测域为 $N\in[0.070542,11.965825]$ B、$D\in[10,299.893]$ B、
$q\in[0.1,1]$；配比为 A4 的 512 条观测配方或其可行凸包点。
主展示用 13 个目标等权和 direct+near 质量政策。在 $N=1$ B、
$D=100$ B 的 512 条观测配方中，108 条满足质量政策；其中
条件 Loss 最小的配方有 $Q_A=1.675449$、$q=0.671361$、
$\widehat L=2.344348$。这是可行观测配方枚举最优，并非连续
凸包或联合预算下的全局最优。

\subsection{参数检验与验证}

表\ref{tab:cyj-q2-validation} 分开列出不同来源真正能检验的内容。
B3 的 8 条插值轨迹对 B1 骨架的平均归一化 RMSE 为 $0.001892$，
B2 的 7 条半合成轨迹为 $0.004370$；二者仅检查形状。
B4/B5 共同支持内的同族规模支配配对，下降方向分别一致 12/12、
7/7；支持外分别是 119/121、106/106，不说明绝对 Loss 可比。

\begin{table}[htbp]
\centering\small
\caption{冻结验证结果及证据等级；B2/B3 为归一化轨迹 RMSE，B4/B5 为方向配对}
\label{tab:cyj-q2-validation}
\begin{tabular}{p{0.24\textwidth}p{0.30\textwidth}p{0.34\textwidth}}
\hline
来源 & 数值 & 解释边界\\ \hline
B1 同源拟合 & RMSE $1.4658\times10^{-4}$ & 同源重构\\
B7 嵌套留 $N/D/Q$ & $0.048738/0.048751/0.048639$ & 半合成同源条件泛化\\
B2/B3 轨迹 & $0.004370/0.001892$ & 半合成或插值，不是独立真实外测\\
B4/B5 支持内方向 & $12/12$、$7/7$ & 同族描述性方向\\
B9/B10 外推压力 & B10 共 128 条，$N$ 全超 B1 上界 & estimated Loss 不作真实测试\\ \hline
\end{tabular}
\end{table}

正式输出及哈希见 CYJ 的 v8 结果清单；发布主体提交
\texttt{615c078} 已经远端核验。18 组冻结接口 fixture
与 95 项 CYJ 回归测试通过，仅说明工程实现与接口一致。
问题三的 CHM 消费者验收仍待完成。

\subsection{边际效用与弹性分析}

设 $q=g(Q_A(\boldsymbol p))$、$F=\exp\{r_w(\boldsymbol p)\}$。
固定配比时，主模型的局部偏导为
\begin{align}\label{eq:cyj-derivatives}
L_N&=F[-\alpha AN^{-\alpha-1}+(1-q)G_N/N],\\
L_D&=F[-\beta BD^{-\beta-1}+(1-q)G_D/D],\\
L_q&=-FG(N,D),\qquad L_{Q_A}=L_qg'(Q_A).
\end{align}
代理截断处只取单侧导数；改变配比时还要计入 $r_w$ 的梯度。
定义无量纲正改善弹性 $\varepsilon_x=-xL_x/L$。在指定点使
$x$ 增加 1\%，Loss 局部相对下降近似为 $\varepsilon_x$\%。

\begin{table}[htbp]
\centering\small
\caption{冻结示例配比、$D=100$ B、$q=0.671361$ 的条件模型局部弹性}
\label{tab:cyj-q2-elasticity}
\begin{tabular}{rrrrr}
\hline
$N$（B 参数） & $\widehat L$ & $\varepsilon_N$ & $\varepsilon_D$ & $\varepsilon_q$\\ \hline
0.1 & 2.780596 & 0.095267 & 0.033814 & 0.115662\\
1 & 2.344348 & 0.055998 & 0.040106 & 0.099511\\
10 & 2.121466 & 0.033091 & 0.044320 & 0.068334\\ \hline
\end{tabular}
\end{table}

本例中参数规模的相对边际收益随 $N$ 增加而下降；该比较固定了
配比和质量代理，且未把不同资源的单位成本视为相同。
若题面成本 $C$ 的质量成本族与基准已指定，同一点的每单位算力
局部收益分别为 $-L_N/C_N$ 与 $-L_q/C_q$；前者大时在该点增加
参数更划算，后者大时提高代理质量更划算。$C_q$ 在质量成本
折点须取可行单侧导数。不同成本族给出的比较可能不同，且
$q$ 由配比派生；故弹性本身不能给出跨情景的唯一支出建议。

\subsection{质量与模型规模的替代关系}

固定 $D,\boldsymbol p$ 作假想的独立质量代理扰动，等 Loss 曲线满足
\begin{equation}\label{eq:cyj-tradeoff}
\frac{dN}{dq}=-\frac{L_q}{L_N},\qquad
\frac{dD}{dq}=-\frac{L_q}{L_D}.
\end{equation}
二者为负表示质量提高后可缩小规模。有限变化还需在支持域内求根。
表\ref{tab:cyj-q2-tradeoff} 区分“旧质量下的等价扩模”和“新质量下
保持原 Loss 的缩模”。

\begin{table}[htbp]
\centering\small
\caption{固定示例配比、$D=100$ B、$q=0.671361$ 的质量--规模条件替代；$N$ 单位 B 参数，根仅在共同支持域内报告}
\label{tab:cyj-q2-tradeoff}
\begin{tabular}{rrrrr}
\hline
初始 $N$ & $dN/dq$ & $q+0.05$ 等价旧质量 $N'$ & $q+0.10$ 等价旧质量 $N'$ & $q+0.05$ 等 Loss 新 $N'$\\ \hline
0.1 & $-0.181$ & 0.109607 & 0.120466 & 0.091384\\
1 & $-2.647$ & 1.144490 & 1.317062 & 0.875789\\
10 & $-30.759$ & 11.697632 & 无支持域内根 & 8.546552\\ \hline
\end{tabular}
\end{table}

例如 $N=1$ B 时，代理质量增加 $0.05$ 的模型内部改善等价于
旧质量下扩模至 $1.144490$ B；若提高质量并保持 Loss 不变，
$N$ 可降至 $0.875789$ B。固定配比却改变由它派生的代理质量，
这里只是敏感性运算，不是已观测的配比干预。

\subsection{领域之间的替代与互补}

配比位于 $\sum_i p_i=1$ 的单纯形内，故只有可行的零和转移方向有
操作含义。从域 $j$ 向域 $i$ 转移 $\epsilon$，在 512 条 A4 配方
的均值点计算
\begin{equation}\label{eq:cyj-pair}
\left.\frac{dL}{d\epsilon}\right|_0=L_{p_i}-L_{p_j},\quad
\left.\frac{d^2L}{d\epsilon^2}\right|_0=
(\boldsymbol e_i-\boldsymbol e_j)^\top H_pL
(\boldsymbol e_i-\boldsymbol e_j).
\end{equation}
梯度同时包括质量代理和配比相对修正两条路径。全部
$\binom{17}{2}=136$ 对均有局部方向和凸包线性规划给出的
可行 $\epsilon$ 上限。例如 freelaw$\to$arxiv 的导数为 $-0.444104$，
该方向上限为 $0.096825$，方向曲率为 $2.816224$；
wikipedia\_en$\to$arxiv 的导数为 $-0.769184$、上限为
$0.075078$。这些方向与互补曲率仅为条件模型内部的局部描述，
没有因果含义。

\begin{table}[htbp]
\centering\small
\caption{A4 观测配方均值点的代表性领域转移；导数为每单位配比转移的条件 Loss，$\epsilon_{\max}$ 来自 A4 凸包}
\label{tab:cyj-q2-pairs}
\begin{tabular}{lrrr}
\hline
转移方向 & $dL/d\epsilon$ & $\epsilon_{\max}$ & 方向曲率\\ \hline
freelaw$\to$arxiv & $-0.444104$ & 0.096825 & 2.816224\\
wikipedia\_en$\to$arxiv & $-0.769184$ & 0.075078 & 0.124064\\ \hline
\end{tabular}
\end{table}

\subsection{外推分析、适用范围与模型局限}

B7 为半合成网格，B3 为 B1 风格插值，B4/B5 未验证统一绝对
Loss 坐标，B10 的 128 条 Loss 为 estimated。其中 B10 的 $N$
全部超过 B1 上界，$D$ 有 $79.6875\%$ 超过上界；不能把这些点
算作真实外推验证。$Q_A\to Q_B$、A 侧逐目标相对 Loss 到 B 侧
总体 Loss 的桥接，以及二者对配比的重复影响，均无法由现有附件
识别。支持域外不直接宣称预测可靠；支持域内的高拟合精度与代码
测试也不证明跨来源科学解释成立。

\subsection{问题二小结}

以 B1 五参数骨架及 B7 三参数半合成质量扩展建立主模型，再以
明示的质量代理与相对配比假设接入问题一。模型在共同支持域内
给出 $N,D,Q,\boldsymbol p$ 的条件局部效用、质量--规模替代及
136 对领域的可行局部方向；跨来源幅度仍待成对实验标定。
问题三接口已发布，CHM 的正式消费者验收另行记录。
